% Part: first-order-logic
% Chapter: sequent-calculus
% Section: derivations

\documentclass[../../../include/open-logic-section]{subfiles}

\begin{document}

\iftag{FOL}
      {\olfileid{fol}{seq}{der}}
      {\olfileid{pl}{seq}{der}}

\olsection{\usetoken{P}{derivation}}

\begin{explain}
প্রারম্ভিক সিকোয়েন্ট দেখতে কেমন এবং অনুমান-বিধিগুলি কী, তা আমরা বলেছি।
সেগুলি থেকেই সিকোয়েন্ট কলনের !!^{derivation}s আরোহীভাবে উৎপন্ন হয়:
প্রতিটি !!{derivation} হয় একটিমাত্র প্রারম্ভিক সিকোয়েন্ট, নয়তো এক বা
দুটি !!{derivation}s-এর পরে বসানো একটি অনুমান নিয়ে গঠিত।
\end{explain}

\begin{defn}[$\Log{LK}$-!!{derivation}]
কোনো সিকোয়েন্ট~$S$-এর একটি \emph{$\Log{LK}$-!!{derivation}} হলো
সিকোয়েন্টের এমন একটি সসীম বৃক্ষ, যা নিচের শর্তগুলি পূরণ করে:
\begin{enumerate}
\item বৃক্ষটির সর্বোচ্চ সিকোয়েন্টগুলি প্রারম্ভিক সিকোয়েন্ট।
\item বৃক্ষটির সর্বনিম্ন সিকোয়েন্ট হলো~$S$।
\item বৃক্ষের $S$ ছাড়া প্রতিটি সিকোয়েন্ট কোনো অনুমান-বিধির সঠিক
  প্রয়োগের পূর্বধারণা, এবং সেই প্রয়োগের সিদ্ধান্তটি বৃক্ষে ওই
  সিকোয়েন্টের ঠিক নিচে থাকে।
\end{enumerate}
তখন $S$-কে !!{derivation}-টির \emph{অন্তিম সিকোয়েন্ট} বলা হয় এবং বলা
হয় যে $S$ \emph{$\Log{LK}$-তে !!{derivable}} (বা
$\Log{LK}$-!!{derivable})।
\end{defn}

\begin{ex}
প্রতিটি প্রারম্ভিক সিকোয়েন্ট, যেমন $!C \Sequent !C$, একটি
!!a{derivation}। এতে, ধরা যাক, $\LeftR{\Weakening}$ বিধি প্রয়োগ করে
নতুন একটি !!{derivation} পাওয়া যায়:
\begin{prooftree}
\Axiom$ \Gamma \fCenter \Delta $
\RightLabel{\LeftR{\Weakening}}
\UnaryInf$ !A, \Gamma \fCenter \Delta$
\end{prooftree}
তবে বিধিটি সাধারণভাবে প্রযোজ্য: বিধির $!A$-কে যেকোনো !!{sentence}
দিয়ে প্রতিস্থাপন করা যায়, যেমন~$!D$ দিয়েও। পূর্বধারণাটি যদি আমাদের
প্রারম্ভিক সিকোয়েন্ট $!C \Sequent !C$-এর সঙ্গে মেলে, তবে $\Gamma$ ও
$\Delta$ উভয়ই শুধু~$!C$, এবং সিদ্ধান্তটি হবে $!D, !C \Sequent !C$।
অতএব নিচেরটি একটি !!a{derivation}:
\begin{prooftree}
\Axiom$ !C \fCenter !C $
\RightLabel{\LeftR{\Weakening}}
\UnaryInf$ !D, !C \fCenter !C$
\end{prooftree}
এবার আরেকটি বিধি, ধরা যাক $\LeftR{\Exchange}$, প্রয়োগ করা যায়; এই
বিধি বাঁদিকে থাকা দুটি !!{sentence}s অদলবদল করতে দেয়। ফলে নিচেরটিও একটি
সঠিক !!{derivation}:
\begin{prooftree}
\Axiom$ !C \fCenter !C $
\RightLabel{\LeftR{\Weakening}}
\UnaryInf$ !D, !C \fCenter !C$
\RightLabel{\LeftR{\Exchange}}
\UnaryInf$ !C, !D \fCenter !C$
\end{prooftree}
এখানে যে বিধিটি প্রয়োগ করা হয়েছে, সেটি দেওয়া হয়েছিল এইভাবে:
\begin{prooftree}
\Axiom$ \Gamma, !A, !B, \Pi \fCenter \Delta $
\RightLabel{\LeftR{\Exchange}}
\UnaryInf$ \Gamma, !B, !A, \Pi \fCenter \Delta,$
\end{prooftree}
এই প্রয়োগে $\Gamma$ ও $\Pi$ উভয়ই ফাঁকা, $\Delta$ হলো $!C$, আর
$!A$ ও $!B$-এর ভূমিকা যথাক্রমে $!D$ ও~$!C$ পালন করেছে। প্রায় একইভাবে
আমরা আরও দেখি যে
\begin{prooftree}
\Axiom$ !D \fCenter !D $
\RightLabel{\LeftR{\Weakening}}
\UnaryInf$ !C, !D \fCenter !D$
\end{prooftree}
একটি !!a{derivation}। এবার এই দুটি !!{derivation}s-কে
$\RightR{\land}$ ব্যবহার করে একত্র করা যায়। বিধিটি ছিল
\begin{prooftree}
\Axiom$\Gamma \fCenter \Delta, !A$
\Axiom$ \Gamma \fCenter \Delta, !B$
\RightLabel{\RightR{\land}}
\BinaryInf$ \Gamma \fCenter \Delta, !A \land !B $
\end{prooftree}
আমাদের ক্ষেত্রে পূর্বধারণাদুটিকে তাদের উপর শেষ হওয়া
!!{derivation}s-এর শেষ সিকোয়েন্টগুলির সঙ্গে মিলতে হবে। অর্থাৎ $\Gamma$
হলো $!C, !D$, $\Delta$ ফাঁকা, $!A$ হলো $!C$ এবং $!B$ হলো $!D$। তাই
অনুমানটি সঠিক হতে হলে সিদ্ধান্ত হবে $!C, !D \Sequent !C
\land !D$।
\begin{prooftree}
\Axiom$ !C \fCenter !C $
\RightLabel{\LeftR{\Weakening}}
\UnaryInf$ !D, !C \fCenter !C$
\RightLabel{\LeftR{\Exchange}}
\UnaryInf$ !C, !D \fCenter !C$
\Axiom$ !D \fCenter !D $
\RightLabel{\LeftR{\Weakening}}
\UnaryInf$ !C, !D \fCenter !D$
\RightLabel{\RightR{\land}}
\BinaryInf$ !C, !D \fCenter !C \land !D $
\end{prooftree}
অবশ্য পূর্বধারণাদুটির ক্রম উল্টেও দেওয়া যায়; তখন $!A$ হবে $!D$ এবং
$!B$ হবে~$!C$।
\begin{prooftree}
\Axiom$ !D \fCenter !D $
\RightLabel{\LeftR{\Weakening}}
\UnaryInf$ !C, !D \fCenter !D$
\Axiom$ !C \fCenter !C $
\RightLabel{\LeftR{\Weakening}}
\UnaryInf$ !D, !C \fCenter !C$
\RightLabel{\LeftR{\Exchange}}
\UnaryInf$ !C, !D \fCenter !C$
\RightLabel{\RightR{\land}}
\BinaryInf$ !C, !D \fCenter !D \land !C $
\end{prooftree}
\end{ex}

\end{document}
